%% workflow-agent.tex — the workflow/agent distinction (Chapter 2). House conceptual style.
%% Two block nodes on a determinism-to-autonomy spectrum, each with its distinctions listed below.
\input{figures/concept-style}
\begin{tikzpicture}
% the spectrum the two sit on
\draw[conceptflow, <->] (-2.1,1.2) -- (8,1.2);
\node[conceptlabel, anchor=south] at (-1.4,1.25) {code specifies each step};
\node[conceptlabel, anchor=south] at (7.4,1.25) {model sets each step};
% the two blocks
\node[conceptbox, fill=blue!15,  minimum width=3cm] (wf) at (0,0) {Workflow};
\node[conceptbox, fill=green!15, minimum width=3cm] (ag) at (6,0) {Agent};
% distinctions, listed below each block (outside the node)
\node[conceptlabel, text width=4.4cm, anchor=north] at (0,-0.75)
  {steps fixed in code\\control flow known in advance\\predictable and cheaper\\chaining, routing, parallelization};
\node[conceptlabel, text width=4.4cm, anchor=north] at (6,-0.75)
  {model chooses next step at run time\\loops until task is completed\\flexible and open-ended\\higher cost, more room for error};
\end{tikzpicture}
